\section{Two further applications: the replicability spectrum}
\label{sec:applications}

The South African audit is the framework's deep case, but a method whose
value is generality must be shown doing different jobs under different data
regimes. This section applies it twice more, at the two ends of the public
data spectrum. In Kenya, where the machine-readable supply--use detail
behind the country's leading SAM is not published, the framework
establishes what \emph{can} be verified from public aggregates---and shows
exactly where verification must stop. In the Netherlands, where complete
harmonised tables are a single API call away, the framework runs in the
opposite direction: not auditing an existing SAM but \emph{generating} one,
with no benchmark consulted at any point. Together with South Africa the
three cases trace a replicability spectrum that we argue is the right way
to think about SAM infrastructure worldwide.

\subsection{Kenya: auditing at the aggregate boundary}
\label{sec:applications:kenya}

The most widely used Kenyan SAM is IFPRI's 2019 Nexus matrix
\citep{thurlow2021}, distributed as an open CC-BY dataset. Its
construction, like that of the Nexus programme generally, reconciles
scattered sources by cross-entropy estimation \citep{robinson2001}, so
cell-level replication of the kind achieved for South Africa is not the
relevant test: a CE-reconciled matrix is a model output, and the honest
question is what its public record allows an outsider to verify. The
framework's answer has three layers.

\emph{Structure verifies.} The distributed file parses into 108 accounts
whose groups (42 activities, 42 commodities, five factors, ten
representative households, and the transaction-cost, enterprise,
government, tax, accumulation, and external accounts) match the data
paper's stated design exactly. The matrix balances only to the file's own
integer rounding---values are printed in whole Sh billions, and account
imbalances reach Sh~5~billion---which is itself a small finding about
distribution practice: rounding a balanced matrix unbalances it.

\emph{Macro controls verify approximately, and the deviations are
informative.} Table~\ref{tab:kenya} compares the SAM's implied national
aggregates with the published current-price tables of the KNBS Economic
Survey 2020. The production side sits a uniform 1.0\,per\,cent below the
published figures---GVA, product taxes, and GDP all shifted
together---which is the signature of a vintage difference in the
then-provisional 2019 accounts, not of a methodological choice
(Section~\ref{sec:discussion}'s class~3). Fixed capital formation,
inventories, and imports agree essentially exactly. The large deviations
are concentrated in private consumption ($-9.2$\,per\,cent), government
consumption ($+9.1$\,per\,cent), and exports ($+13.6$\,per\,cent), and
they have a structural explanation: KNBS's own expenditure table carries a
production-versus-expenditure discrepancy of Sh~$-352$~billion
(3.6\,per\,cent of GDP), and a balanced SAM cannot carry a statistical
discrepancy---the reconciliation \emph{must} allocate it somewhere. The
deviation pattern is consistent with the discrepancy, together with
associated source and definitional differences, being absorbed primarily
through the consumption and trade aggregates; aggregate comparisons
cannot identify the allocation uniquely, and that very limit is the
point. Nothing in the public record documents the allocation at the
level where it acts, which is precisely the provenance gap the paper's
proposal addresses.

\begin{table}[t]
\centering\small
\caption{Kenya: the IFPRI 2019 Nexus SAM against published KNBS controls
(Sh billion, current prices; Economic Survey 2020, Tables 2.1 and 2.7,
2019 provisional).}
\label{tab:kenya}
\begin{tabular}{lrrr}
\toprule
Aggregate & SAM & KNBS & Dev.\ \% \\
\midrule
GVA (all activities) & 8,825 & 8,912 & -1.0 \\
Taxes on products & 821 & 829 & -0.9 \\
GDP at market prices & 9,646 & 9,740 & -1.0 \\
Private consumption (incl. NPISH) & 7,295 & 8,035 & -9.2 \\
Government consumption & 1,387 & 1,271 & +9.1 \\
Gross fixed capital formation & 1,631 & 1,632 & -0.1 \\
Changes in inventories & 63 & 63 & -0.2 \\
Exports of goods and services & 1,331 & 1,172 & +13.6 \\
Imports of goods and services & 2,071 & 2,082 & -0.5 \\
Statistical discrepancy & 0 & -352 & --- \\
\bottomrule
\end{tabular}

\end{table}

\emph{Cell-level audit is bounded by publication practice, not by
method.} The SAM's supply--use core descends from KNBS's 2009 benchmark
tables, whose full balanced database (151 products by 81 industries) is,
per the Economic Survey, available from KNBS on request but published only
as condensed print tables. We did not request it, and the boundary tested
here is therefore a specific one: what an auditor can verify from the
public record alone, without privileged access or correspondence with the
producer. That is the boundary a reader of a published SAM faces, and the
one the paper's standards proposal addresses; an audit conducted on
requested files would test a different and more permissive question. Two
consequences follow. The Kenyan evidence below is confined to structure,
grain, and macro controls---the depth the public record supports---and the
claim that the South African pipeline would extend to the Kenyan detail
remains untested: the reader for the KNBS layout, and its concordances,
would still have to be written, even though the method behind them is
unchanged. That boundary---audit feasible in principle, unreachable from
what is published---defines the middle of the replicability spectrum.

\subsection{The Netherlands: benchmark-free generation}
\label{sec:applications:nl}

The third application inverts the exercise. For a country inside the
European statistical system, the framework should need no benchmark, no
technical note, and no manual acquisition: the ESA transmission programme
puts harmonised supply--use tables and sector accounts on a public API
\citep{eurostat2008}. We test that claim on the Netherlands for 2019.
Three API calls retrieve the supply table, the use table, and the full
sequence of non-financial sector accounts (about 180~KB in total); a
70-line JSON-stat reader added to the toolkit decodes them, and a single
script assembles a macro accounting matrix and balances it into a macro
SAM. No published SAM was consulted at any point during construction or
validation---though, as the results below make clear, that is a statement
about SAMs and not a claim of independent verification: the checks
available here are internal consistency and agreement with published
national-accounts aggregates drawn from the same compilation system that
supplies the inputs.

Two results follow. First, \emph{the official tables are internally
exact}: every identity the toolkit enforces---output and value-added
identities for all 65 industries, the import identity, and the
supply--demand balance of all 65 products---closes at a EUR~1~million
tolerance with zero findings. The contrast with the audit cases is the
point: when a statistical office publishes machine-readable tables under a
harmonised standard, deterministic verification is free. Even here,
however, the framework's refusal to classify by convention earned its
keep: Eurostat's category lists mix aggregates with overlapping
sub-details (NACE section codes alongside their own divisions), and a
na\"ive selection double-counts; the toolkit's coverage-based partition,
guarded by an assertion that the chosen codes tile the total exactly,
caught both this and a code collision between the NACE section~P
(education) and the ESA transaction P-codes. Names mislead in Luxembourg
as in Pretoria.

Second, \emph{the generated matrix closes to within 1.5\,per\,cent of
GDP with every residual attributed---and a logged balancing step then
closes it exactly}. We are careful with terminology here: a SAM is
defined by every account balancing, so the direct output of the
generation---the 18-account \emph{macro accounting matrix}---is the
diagnostic product, and the SAM is what balancing makes of it. The matrix
routes factor income, taxes, transfers, and the ESA bridging items (the
pension-fund adjustment D8, capital transfers D9, net lending B9, and the
residents-abroad consumption adjustments) through explicit instrument
accounts; GDP at basic prices emerges at EUR~724{,}960~million, on the
published figure, and compensation of employees closes to the euro across
four series. Both agreements are internal to the Dutch national accounts:
the four series and the published aggregate are outputs of one compilation
system, so what they establish is that our assembly preserves CBS's own
consistency, not that it is independently corroborated. Table~\ref{tab:netherlands} reports the residual
of every account before balancing: twelve are exactly zero, and the six
that are not are all below 1.5\,per\,cent of GDP and attach to known ESA
boundary items (the domestic-versus-national concept in compensation,
mixed-income attribution, capital-account detail). Consistent with the
framework's first commitment, the residuals are published, not hidden;
the balanced macro SAM is then produced by the toolkit's logged RAS
step: it converges to a maximum residual below EUR~0.01~million at full
precision, with all 70 adjustment factors inside $[0.9832, 1.0282]$,
concentrated on the rest-of-world cells---exactly where the attributed
residuals sat---and every factor ships with the matrix. One caveat
applies to the distributed artefact, and it is the same one this paper
notes for the Kenyan file: the CSV prints values at EUR~0.1~million
precision, so recomputing account residuals from the distributed file
yields rounding imbalances of up to EUR~0.2~million; the full-precision
matrix regenerates from the replication package. The balancing choices (the treatment
of two negative cells, the target vector, and the choice of RAS), the
invariants they preserve, and a target-vector sensitivity check are
documented in Appendix~\ref{app:nlbalance}.

\begin{table}[t]
\centering\small
\caption{Netherlands: account residuals of the generated 2019 macro
accounting matrix, before balancing (no benchmark consulted; GDP at basic
prices EUR 724{,}960m). The logged RAS step closes all residuals with
adjustment factors in $[0.9832, 1.0282]$.}
\label{tab:netherlands}
\begin{tabular}{lrr}
\toprule
Account & Residual (EUR m) & \% of GDP \\
\midrule
Corporations & +1,532 & +0.21 \\
Government & +1,443 & +0.20 \\
Households & +9,111 & +1.26 \\
Labour & -4,193 & -0.58 \\
Rest of world & -10,830 & -1.49 \\
Savings--investment & +2,937 & +0.41 \\
\midrule
\multicolumn{3}{l}{\emph{Remaining 12 accounts (production, taxes, capital, instruments): residual zero}} \\
\bottomrule
\end{tabular}

\end{table}

The cost ledger deserves its own sentence, because it carries the paper's
economic claim. The South African reconstruction required four public
providers, a registered microdata access, and bulk downloads diagnosed
around a broken online query system; the Kenyan case required locating
condensed print tables inside a 322-page yearbook; the Dutch case required
three API calls and a single working session, and the 70-line reader that
enabled it works unchanged for every country in the ESA transmission
programme. Out-of-sample runs for Austria and Poland confirm that the
generator transfers---and that confidentiality suppression remains a
cell-level constraint even under harmonised publication;
Appendix~\ref{app:transfer} reports the results. Conditional on the
toolkit, the marginal implementation effort for a verified macro SAM is
therefore shaped largely by the accessibility and structure of the
published data---the replicability spectrum restated as effort.
